\documentclass[11pt]{article}

\usepackage[utf8]{inputenc}
\usepackage[T1]{fontenc}
\usepackage[margin=1in]{geometry}
\usepackage{amsmath,amssymb,amsthm,mathtools}
\usepackage{hyperref}
\usepackage{xcolor}
\usepackage{enumitem}
\usepackage{setspace}
\usepackage{csquotes}

\hypersetup{
  colorlinks=true,
  linkcolor=blue,
  citecolor=blue,
  urlcolor=blue
}

\onehalfspacing

\newtheorem{definition}{Definition}
\newtheorem{assumption}{Assumption}
\newtheorem{proposition}{Proposition}
\newtheorem{theorem}{Theorem}
\newtheorem{remark}{Remark}

\title{A Formal Model of an Open Technological Skill-Development Ecosystem}
\author{
  Author Name \\
  Universidad Autonoma Metropolitana Cuajimalpa \\
  \texttt{author@email.com}
}
\date{\today}

\begin{document}

\maketitle

\begin{abstract}
This paper introduces the foundational structure of a research program on open technological skill-development ecosystems. Its purpose is to define the object of study, state the main research problem, identify the principal design objectives, and prepare the conceptual-formal transition required for later axioms, definitions, propositions, and degradation analyses.
\end{abstract}

\noindent\textbf{Keywords:}
skill-development ecosystems, socio-technical systems, mechanism design, applied mathematics, sustainable competition

\section{Introduction}

\paragraph{Purpose of this section.}
Introduce the general problem, explain why some ecosystems should be studied as environments for capability formation rather than merely as event or transaction systems, and motivate the need for a conceptual-formal treatment.

Some competitive systems are valuable not only because they coordinate events, but because they shape the long-term development of participants. In such cases, the central object of study is not the event itself, but the ecosystem that makes repeated participation, learning, and progression possible.

This paper is motivated by the idea that technological competition can be designed as a developmental environment. If that is true, then the relevant question is not only how to operate the system, but how to preserve the structural conditions under which capability growth, accessibility, and economic continuity can coexist.

A central intuition behind this work is that development does not emerge automatically from repeated competition. Repetition may also produce concentration, discouragement, opportunism, or exclusion. For that reason, the developmental character of the ecosystem must be treated as a design problem rather than as a byproduct of participation.

The paper therefore starts from a shift in emphasis: the relevant unit of analysis is neither the isolated participant nor the isolated event, but the evolving structure that links entry, repetition, incentives, technological experimentation, and future opportunity.

\subsection{Why this paper exists}

\subsection{Why the problem matters}

\subsection{Why the ecosystem must be studied formally}

\section{Research Problem}

\paragraph{Purpose of this section.}
State the central research problem in natural language and identify the main tensions that the ecosystem must resolve.

The research problem is to understand how an open technological ecosystem can remain developmental without becoming exclusionary, opportunistic, or economically fragile. The difficulty lies in the fact that openness, competition, innovation, and sustainability do not automatically reinforce one another.

The system must therefore be designed so that its incentives and constraints preserve the conditions for skill development over time. This requires identifying which structural properties are essential and which forms of degradation arise when they are weakened.

At least four tensions appear immediately. First, lowering barriers to entry may increase heterogeneity and operational complexity. Second, rewarding strong performance may unintentionally magnify prior economic asymmetries. Third, allowing economic incentives may energize participation while also making strategic abuse more attractive. Fourth, technological freedom may stimulate innovation while simultaneously raising the cost of staying competitive.

These tensions suggest that the problem is not merely one of coordination. It is one of balance under repeated interaction. The ecosystem must be open enough to remain developmental, strict enough to remain credible, and economically structured enough to remain alive.

\subsection{General statement of the problem}

\subsection{Main tensions}

\subsection{Research question}

\subsection{Subquestions}

\section{Intuitive View of the Ecosystem}

\paragraph{Purpose of this section.}
Describe the ecosystem at a high level before any heavy formalization. Clarify what kind of object is being studied and what distinguishes it from a league, a platform, or a simple event-management system.

The ecosystem studied here is not best understood as a sequence of isolated competitions. It is better understood as a structured environment in which agents repeatedly interact, acquire capabilities, test technologies, and build trajectories of participation.

Its identity depends on continuity, progression, and role diversity. The competitive layer matters, but it matters as part of a broader developmental system rather than as an end in itself.

This means that the ecosystem must be able to connect present participation with future possibility. A newcomer should not merely be allowed to appear once; the system should make it meaningful to continue. A technical improvement should not merely affect one event; it should become part of a pathway of capability formation. A role such as participant, promoter, owner, or organizer should not remain isolated; it should become part of a larger structure of developmental mobility.

The technological nature of the ecosystem is also important. Skill is not restricted here to athletic execution. It includes design, adaptation, resource allocation, technical judgment, and strategic decision-making. For that reason, the ecosystem should be able to support both human performance and technological learning within the same developmental frame.

\subsection{Why this is an ecosystem}

\subsection{Why this is about skill development}

\subsection{Why the ecosystem is technological}

\subsection{Why openness matters}

\section{Design Objectives}

\paragraph{Purpose of this section.}
List and organize the system-level objectives that the ecosystem is intended to preserve.

The ecosystem should not be optimized along a single axis. A system that maximizes performance while destroying accessibility is not satisfactory, and a system that maximizes openness while tolerating opportunism is also unsatisfactory.

The design objectives must therefore be treated as a constrained bundle. Accessibility, development, sustainability, enforceability, and innovation compatibility are not independent goals, but jointly necessary properties of the ecosystem.

Their relationship is not symmetric. Some objectives constrain others. Accessibility constrains how much advantage prior resources may legitimately generate. Enforceability constrains how much informality the system can tolerate. Sustainability constrains how much idealized openness can be supported without bounded commitments or recurrent support structures. Innovation compatibility constrains how rigidly the system may regulate technical experimentation.

As a result, the design objectives should not be interpreted as a checklist. They form an interdependent system of requirements. If one objective is pushed too aggressively while the others are ignored, the ecosystem may survive operationally while ceasing to be developmental in the intended sense.

\subsection{Accessibility}

\subsection{Developmentality}

\subsection{Economic sustainability}

\subsection{Enforceability}

\subsection{Innovation compatibility}

\section{Informal Axioms}

\paragraph{Purpose of this section.}
State the first structural principles of the ecosystem in plain language, before translating them into formal assumptions.

If the ecosystem is to preserve its developmental character, then some design choices cannot be treated as optional preferences. They must be regarded as structural conditions.

These conditions are introduced first in informal language because their role must be conceptually clear before they are formalized. Their function is to express what the ecosystem must preserve in order not to degrade into a different kind of system.

The reason to speak of axioms here is not that the system is already fully formal, but that certain commitments appear prior to implementation detail. For example, if future opportunity is not linked to present participation, the ecosystem may still organize events, but it no longer functions as a developmental structure. If violations carry no durable consequence, the ecosystem may still be open, but it no longer preserves trust under repeated interaction.

Informal axioms therefore serve as pre-formal commitments. They identify the conceptual backbone that later definitions and propositions will attempt to express with greater precision.

\subsection{Entry must remain open}

\subsection{Participation must create future opportunity}

\subsection{Advantages must remain bounded}

\subsection{Commitments must have consequences}

\subsection{Economic promises must remain bounded}

\subsection{Multiple roles must be possible}

\section{Formal Model}

\paragraph{Purpose of this section.}
Introduce the first formal objects of the paper: agents, roles, variables, state descriptions, and system-level quantities.

The formal model is introduced only after the main intuition has been established. Its purpose is not to replace the conceptual view, but to make it analyzable.

At this stage, the goal is to identify the minimum set of formal objects required to describe development, access, incentives, and degradation. Later versions of the paper will strengthen these objects and relate them more explicitly.

The initial model should remain deliberately minimal. If too many variables are introduced too early, the formal layer may become locally sophisticated but globally uninformative. The first task is instead to mark the main dimensions along which the ecosystem evolves: capability, resources, reputation, participation, and system-level viability.

This also means that not every important phenomenon needs immediate formal detail. Some objects may first appear as placeholders whose exact mathematical role will be refined in later layers. What matters in the present stage is to establish the analytical space within which stronger claims will later be made.

\subsection{Agents}

\subsection{Roles}

\subsection{Participant-level variables}

\subsection{Event-level variables}

\subsection{Ecosystem-level variables}

\subsection{System objective}

\section{Structural Properties}

\paragraph{Purpose of this section.}
Identify the first properties that should be preserved by the system and prepare the ground for later propositions and theorems.

The ecosystem is meaningful only if some of its intended properties remain stable under repeated interaction. These properties are not merely descriptive; they are the reason the ecosystem is worth designing in the first place.

The main task of this section is to identify those properties and prepare the transition from informal design objectives to formal claims about what the system can or cannot preserve.

In this sense, structural properties play the role of invariants or quasi-invariants of the ecosystem. They need not be absolute in every local situation, but they must remain sufficiently preserved at the system level if the ecosystem is to retain its identity.

Examples of such properties include meaningful entry, non-trivial progression, bounded domination by prior wealth, and non-negligible consequences for non-compliance. Later versions of the paper should ask not only whether the system exhibits these properties under ideal conditions, but also how robustly they survive under stress or asymmetry.

\subsection{Properties to preserve}

\subsection{Why those properties depend on structure}

\subsection{First candidate propositions}

\section{Degradation Modes}

\paragraph{Purpose of this section.}
Describe the major ways in which the ecosystem can fail if key structural conditions are removed or weakened.

The ecosystem must be studied not only through its intended goals, but also through its characteristic failure modes. A design becomes more intelligible when one can say not only what it is trying to achieve, but what kinds of collapse or distortion it is explicitly trying to avoid.

Degradation is especially important here because many harmful outcomes are structurally plausible: exclusion, wealth capture, rational opportunism, or loss of developmental meaning. These are not accidental pathologies but candidate outcomes of weak design.

Thinking in terms of degradation also prevents a common mistake: assuming that the absence of immediate failure is equivalent to health. A system may continue to operate while slowly losing openness, trust, or developmental value. For that reason, degradation must be treated as a family of trajectories, not merely as a binary failure event.

Each degradation mode should eventually be linked to one or more weakened structural conditions. This makes the theory useful not only for description but for diagnosis: it becomes possible to explain how a system drifts away from its intended form and which constraints are being lost in the process.

\subsection{Exclusion}

\subsection{Plutocratic drift}

\subsection{Opportunistic drift}

\subsection{Economic fragility}

\subsection{Loss of developmental value}

\section{Discussion}

\paragraph{Purpose of this section.}
Clarify the scope of the present paper, distinguish what is already being claimed from what remains future work, and explain the intended audience and interpretation of the framework.

This paper is intentionally foundational. It does not aim to provide a complete mechanism, a full equilibrium analysis, or an empirical validation of the ecosystem.

Its role is to define the object of study, establish a common language, and identify the first layer of structural commitments needed for later mathematical and empirical work.

Because of this foundational role, the paper should be readable by more than one community. A reader from applied mathematics should be able to see where the formal development is heading. A reader from socio-technical systems should be able to understand the institutional and developmental meaning of the framework. A reader from competition, sport, or technology studies should be able to recognize the practical relevance of the problem.

This multi-audience ambition is not incidental. It follows from the object itself. The ecosystem under study is simultaneously mathematical, institutional, technological, and developmental.

\subsection{What this paper does}

\subsection{What this paper does not yet do}

\subsection{Interpretive scope}

\section{Research Agenda}

\paragraph{Purpose of this section.}
Position the present article as the first contribution in a broader line of research and identify the next papers that naturally follow.

The present article should be read as the opening move of a larger research program. Once the ecosystem is defined conceptually and structurally, one can study specific aspects of it with greater precision.

The most natural continuations concern reputation, bounded incentives, accessibility under resource asymmetry, dynamic evolution, and empirical observation. Each of these topics can become a paper in its own right.

The order of that agenda matters. Reputation and bounded commitments should likely be addressed early because they are closely connected to enforceability and economic credibility. Accessibility under asymmetry should also appear early because it bears directly on the anti-exclusionary identity of the ecosystem. Simulation can then serve as a bridge between conceptual-formal reasoning and domain-level plausibility.

Empirical validation should arrive after the conceptual and computational layers are sufficiently articulated. Otherwise, empirical observation may remain descriptive without testing the structural claims that motivate the framework.

\subsection{Paper on reputation and compliance}

\subsection{Paper on accessibility versus plutocratic drift}

\subsection{Paper on bounded economic promises}

\subsection{Paper on multi-agent simulation}

\subsection{Paper on empirical validation}

\section{Conclusion}

\paragraph{Purpose of this section.}
Close the paper by restating the object of study, the core problem, and the role of the present article as a foundation for future formal and empirical work.

The central claim of this paper is that some competitive systems should be studied as technological skill-development ecosystems rather than as mere operational platforms. This shift changes both the object of analysis and the mathematical questions that become relevant.

The contribution of the paper is therefore foundational: it proposes the object, frames the problem, and prepares the structure within which stronger definitions, propositions, and later validations can be developed.

The broader implication is that the success of such a system cannot be reduced to event execution, participation counts, or short-term visibility. If the ecosystem is truly developmental, then its quality depends on whether it preserves the conditions under which people can enter, learn, progress, innovate, and remain meaningfully connected over time.

The next layer of work should therefore move from conceptual organization to pre-formal commitments and then to the first explicit formal definitions of the ecosystem's core properties.

\bibliographystyle{plain}
\begin{thebibliography}{9}

\bibitem{placeholder1}
Placeholder reference.

\end{thebibliography}

\end{document}
